\documentclass[11pt]{article}
\usepackage[margin=0.9in]{geometry}
\usepackage{amsmath,amssymb,amsthm,mathtools,hyperref}
\hypersetup{hidelinks,pdftitle={Q222: exact canonical completion and polynomial-selector obstruction}}
\newtheorem{theorem}{Theorem}
\newtheorem{proposition}[theorem]{Proposition}
\newtheorem{corollary}[theorem]{Corollary}
\newcommand{\PSD}{\succeq0}
\newcommand{\R}{\mathbb R}
\title{Canonical Bell validity through two matrix inequalities}
\author{}
\date{}
\begin{document}
\maketitle
\vspace{-2em}
\begin{abstract}
A Bell inequality must hold for every state and both measurement angles.
For the canonical $Q_{222}$ slack, a polynomial change of state coordinates
eliminates one angle: validity is equivalent to two $4\times4$ matrix
inequalities with a $2\times2$ completion depending on the remaining angle.
The equivalence includes zero source weights and projective endpoints.
We then give a rational valid inequality for which no polynomial completion
exists. A contact fixes its three possible polynomial parameters, and a
two-coordinate quadratic test excludes them. A fixed algebraic completion or a
common testing multiplier would be needed to obtain a finite global
certificate by this construction.
\end{abstract}

\section{Eliminating one measurement angle}
Let $c=(d,\lambda,n,m,C,A,D,\alpha,\beta)\geq0$ and put
$D_u=1+u^2$, $D_v=1+v^2$.  The established complete canonical parent is
\begin{align}
F_c={}&4\lambda D_u(W+vX)^2+4nD_uD_vX^2
 +4m(vW-X+uvY-uZ)^2\notag\\
&+4CD_v(uW-Y)^2+4AD_uD_vY^2
 +(4D+2d)(uW+uvX-Y-vZ)^2\notag\\
&+(4\alpha+2d)D_u(vY-Z)^2
 +(4\beta+2d)D_v(X+uZ)^2-2dD_uD_vZ^2.\label{parent}
\end{align}
Define the symmetric
$2\times2$ matrices
\begin{align}
P(v)&=\begin{pmatrix}
4mv^2&-4mv\\-4mv&4m+(4\beta+2d)D_v
\end{pmatrix},\notag\\
N(v)&=\begin{pmatrix}
4CD_v+4D+2d&(4D+2d)v\\(4D+2d)v&(4D+2d)v^2
\end{pmatrix},\label{blocks}\\
E(v)&=\begin{pmatrix}
4AD_v+(4\alpha+2d)v^2&-(4\alpha+2d)v\\
-(4\alpha+2d)v&4\alpha-2dv^2
\end{pmatrix},\notag\\
K(v)&=4\lambda\begin{pmatrix}1&v\\v&v^2\end{pmatrix}
       +4nD_v\begin{pmatrix}0&0\\0&1\end{pmatrix}.\notag
\end{align}
All entries are affine in $c$ and quadratic in the testing variable $v$.
For $a=(S,T)^T$, $b=(P_0,Q_0)^T$, and $U=u^2$, make the polynomial change
\begin{equation}
W=S+UP_0,\quad X=T+UQ_0,\quad
Y=u(S-P_0),\quad Z=u(T-Q_0).\label{change}
\end{equation}
The state coordinates $P_0,Q_0$ are unrelated to the matrix $P(v)$.
Direct substitution gives the exact identity
\begin{equation}
F_c=D_u\Phi_c(U,v;a,b),\qquad
\boxed{\begin{aligned}
\Phi_c={}&(1+U)a^TPa+U(1+U)b^TNb\\
 &+U(a-b)^TE(a-b)+(a+Ub)^TK(a+Ub).
\end{aligned}}\label{normal}
\end{equation}
For example, the $m,C,D,\alpha,\beta$ event roots in \eqref{parent}
become respectively
$D_u(vS-T)$, $uD_uP_0$, $uD_u(P_0+vQ_0)$,
$u\{v(S-P_0)-(T-Q_0)\}$, and $D_uT$.
These substitutions prove \eqref{normal} without a source-dependent inverse.

Write $\psi=(W,X,Y,Z)^T$ and let $\mathcal P_c(U,v)$
be the symmetric matrix of $\Phi_c$ in $(a,b)$. Its three coefficients are
\begin{equation}
\mathcal P_c=
\begin{pmatrix}P+K&0\\0&0\end{pmatrix}
+U\begin{pmatrix}P+E&K-E\\K-E&N+E\end{pmatrix}
+U^2\begin{pmatrix}0&0\\0&N+K\end{pmatrix}.\label{quadratic}
\end{equation}
The change matrix has determinant $u^2D_u^2$.  Its polynomial inverse
numerator is
\begin{equation}
L(u)=\begin{pmatrix}
u&0&u^2&0\\0&u&0&u^2\\u&0&-1&0\\0&u&0&-1
\end{pmatrix},\qquad
\boxed{u^2D_uF_c(\psi)=(L(u)\psi)^T
\mathcal P_c(u^2,v)(L(u)\psi).}\label{multiplier}
\end{equation}
To turn \eqref{multiplier} into squares, we first factor its quadratic dependence
on $U$.

\section{A pointwise matrix completion}
\begin{theorem}\label{thm:selector}
For any fixed real source $c$ and fixed real Bob test $v$, the following are
equivalent:
\begin{enumerate}
\item $\mathcal P_c(U,v)\PSD$ for every $U\geq0$.
\item There exists a real $2\times2$ matrix $J$ such that
\begin{equation}
\boxed{\mathcal G_J=
\begin{pmatrix}P+K&J\\J^T&N+K\end{pmatrix}\PSD,
\qquad
\mathcal H_J=
\begin{pmatrix}P+E&K-E-J\\K-E-J^T&N+E\end{pmatrix}\PSD.}\label{lmis}
\end{equation}
\end{enumerate}
The equivalence also holds over every real closed field.  In particular,
complete canonical validity is exactly
\begin{equation}
\boxed{\forall v\in\R\ \exists J\in\R^{2\times2}:
\mathcal G_J(v)\PSD,\quad\mathcal H_J(v)\PSD.}\label{quantifiers}
\end{equation}
All source faces, including $d=0$ and arbitrary zero event weights, are retained.
\end{theorem}
\begin{proof}
The reverse implication is the explicit identity
\begin{equation}
\Phi_c=
\binom a{Ub}^{\!T}\mathcal G_J\binom a{Ub}
+U\binom ab^{\!T}\mathcal H_J\binom ab.\label{split}
\end{equation}
For the forward implication, $\mathcal P_c(t^2,v)$ is a positive semidefinite
univariate matrix polynomial of degree at most four.  The real-line matrix
Fej\'er--Riesz theorem gives a polynomial matrix factor of degree at most two;
real and imaginary parts give a real sum of matrix squares
\cite[Theorem 1.2]{Zalar}. Averaging its values at $t$ and $-t$ separates even
and odd powers. Consequently
\[
\mathcal P_c(U,v)=(B_0+UB_2)^T(B_0+UB_2)+UB_1^TB_1
\]
for rectangular real matrices $B_0,B_1,B_2$.
The constant coefficient in \eqref{quadratic} has its last two rows and columns
zero, so $B_0=[R\ \ 0]$.  The coefficient of $U^2$ has its first two rows and
columns zero, so $B_2=[0\ \ T]$.  This uses only the elementary fact that a
zero diagonal entry of a positive semidefinite Gram matrix forces its entire
row and column to vanish.  Setting $J=R^TT$ gives
$\mathcal G_J=[R\ T]^T[R\ T]$ and
$\mathcal H_J=B_1^TB_1$, proving \eqref{lmis}.

This is a fixed finite-dimensional first-order assertion in the real closed
field language: positivity is expressed by principal minors, and the
half-line condition by quantification over $U$ and a four-vector.  Transfer
therefore proves precisely the same assertion over every real closed field.

For $u\ne0$, \eqref{change} is invertible and so positivity of $F_c$ is equivalent
to positivity of $\mathcal P_c(u^2,v)$.  The excluded value $u=0$ is recovered
by polynomial continuity; all $U>0$ are attained.  More explicitly,
\eqref{parent} has bidegree at most $(2,2)$ in $(u,v)$, so
\[
F_c^{\mathrm h}(r,s;\psi)
=r_0^2s_0^2F_c(r_1/r_0,s_1/s_0;\psi)
\]
is polynomial in the homogeneous angle pairs $r=(r_0,r_1)$ and
$s=(s_0,s_1)$.  On their unit circles it equals the normalization
$F_c/(D_uD_v)$.  Continuity recovers the Alice and Bob projective endpoints.
Over a real closed field the same endpoint statement follows by the sign of
the first nonzero coefficient in the corresponding local polynomial chart.
No source coefficient was divided out.
\end{proof}

There is also an elementary check on the face $d=0$: then $E\PSD$, and
$J=K$ makes both matrices in \eqref{lmis} sums of the positive diagonal blocks
$P,N$ and a square block built from $K$ or $E$. More generally adding a
nonnegative event source with $d=0$ preserves any certificate by adding its
$K$ to $J$.  The unresolved completion is therefore associated with the signed
core, not with adding positive events.

The remaining angle enters the completion through $J(v)$. A finite global
certificate would require these matrices to lie in one fixed function space.
Pointwise transfer permits algebraic dependence on both the source and the
angle; it supplies no such space in the tensor product of their independent
fields. We therefore examine the simplest possibility, polynomial $J(v)$.

\section{Why polynomial completions do not suffice}
\begin{proposition}\label{prop:degree}
If polynomial entries $J(v)$ satisfy \eqref{lmis} for every real $v$, they have
degree at most two.  A selector can then be chosen of the form
\[
J(v)=\begin{pmatrix}p_0+p_2v^2&q_1v\\r_1v&s_0+s_2v^2\end{pmatrix}.
\]
On the entire C1 face $\lambda=n=0$, existence of any polynomial selector is
equivalent to existence of one with just three scalar parameters:
\begin{equation}
\boxed{J(v)=\begin{pmatrix}pv^2&0\\rv&sv^2\end{pmatrix}.}\label{three}
\end{equation}
\end{proposition}
\begin{proof}
Every $2\times2$ principal minor of $\mathcal G_J$ gives
$|J_{ij}(v)|^2\leq(P+K)_{ii}(N+K)_{jj}$.  The right side has degree at most
four, proving the degree bound, including cases where it is identically zero.
With $T=\operatorname{diag}(1,-1)$, each matrix in \eqref{blocks} obeys
$M(-v)=TM(v)T$.  Replace $J(v)$ by
$\tfrac12\{J(v)+TJ(-v)T\}$; convexity of the positive semidefinite cone retains
both inequalities and gives the displayed parity.

On C1, $P(0)=\operatorname{diag}(0,4m+4\beta+2d)$ and
$N(0)=\operatorname{diag}(4C+4D+2d,0)$.  Zero diagonal entries in
$\mathcal G_J(0)$ force $p_0=s_0=0$. Finally
\[
q_1^2v^2=|J_{12}(v)|^2\leq4m(4D+2d)v^4
\]
for every $v$, so $q_1=0$. This proof does not assume any of the source weights
is positive.
\end{proof}

A small rational example already excludes polynomial completions. Set
$\lambda=n=0$ and
\[
 (d,m,C,A,D,\alpha,\beta)
 =\left(359,\frac{75}{16},60,416,\frac{41}{2},
                       \frac{831}{2},\frac{151}{2}\right).
\]
Let $P_*(u,v)$ be the four-dimensional state matrix of \eqref{parent}
at this source. Its validity follows from four elementary AM--GM
inequalities. Put $\xi=4u^2$, $\eta=v^2/4$, and
\begin{align*}
 M(\xi,\eta)={}&3757645(\xi+\xi^{-1}-2)
       +14022297(\eta+\eta^{-1}-2)\\
 &+3732435(\eta/\xi+\xi+\eta^{-1}-3)\\
 &+8179670(\xi\eta+\xi^{-1}+\eta^{-1}-3).
\end{align*}
Each bracket is nonnegative for $\xi,\eta>0$: its positive terms have
product one. Direct expansion gives
\begin{equation}
 \det P_*(u,v)=28800D_u^2D_v^2u^2v^2M(\xi,\eta).
 \label{valid-det}
\end{equation}
The $Z=0$ principal block $B_*$ is strictly positive when $u\ne0$.
Indeed, the positive $m,C,A$ terms of \eqref{parent} give
\[
 B_*\succeq T^T\operatorname{diag}(4m,4CD_v,4AD_uD_v)T,
 \qquad
 T=\begin{pmatrix}v&-1&uv\\u&0&-1\\0&0&1\end{pmatrix},
 \quad \det T=u.
\]
All remaining terms on $Z=0$ are nonnegative. The scalar Schur complement
therefore proves $P_*\succeq0$ for $uv\ne0$ from \eqref{valid-det}.
Polynomial continuity supplies the finite axes, and the homogeneous
angle charts supply the projective endpoints. Separating real and imaginary
parts covers complex states as well.

The contact has $u_*=1/2$, $v_*=2$ and
$\psi_*=(2,-1/4,1/4,1)^T$. Under \eqref{change} it becomes
\begin{equation}
 U_*=\frac14,\qquad
 a_*=\binom{17/10}{1/5},\qquad
 b_*=\binom{6/5}{-9/5}.
 \label{contact}
\end{equation}
Substitution verifies the full state equation
$\mathcal P_c(U_*,v_*)(a_*,b_*)^T=0$.

\begin{theorem}\label{thm:obstruction}
This valid source has no polynomial selector satisfying \eqref{lmis} for all
$v$. Equivalently, its polynomial $\Phi_c(u^2,v;a,b)$ is not a sum of polynomial
squares in $(u,v,a,b)$.
\end{theorem}
\begin{proof}
By Proposition~\ref{prop:degree}, a polynomial selector would have the form
\eqref{three}. At the contact, both nonnegative terms in \eqref{split} vanish.
Therefore
\begin{equation}
J(v_*)b_*=-P(v_*)a_*/U_*,\qquad
J(v_*)^Ta_*=-U_*N(v_*)b_*.\label{kernel}
\end{equation}
The first coordinate of the first equation fixes $p$, and the second
coordinate of the second fixes $s$. Either remaining coordinate then gives
$r$. With \eqref{contact}, the unique solution is
\begin{equation}
 p=-100,\qquad r=2000,\qquad s=1200.
 \label{forced}
\end{equation}
At $v=1$, the principal submatrix of $\mathcal G_J$ on
$(a_2,b_1)$ would consequently be
\[
 \begin{pmatrix}
 8235/4&2000\\
 2000&1280
 \end{pmatrix}.
\]
Its quadratic value on $(1,-1)$ is $-2645/4$, contradicting
$\mathcal G_J(1)\PSD$.

Finally, a polynomial sum of squares for $\Phi_c(u^2,v;a,b)$ can use roots
of degree at most two in $u$, at most one in $v$, and homogeneous linear
degree in the four state coordinates: higher terms cannot cancel in a sum
of real squares. Averaging in $u\mapsto-u$ and using the constant and leading
state supports exactly as in Theorem~\ref{thm:selector} would produce a
polynomial selector. Conversely a polynomial selector gives polynomial
matrix factorizations of its two positive univariate matrices by
Fej\'er--Riesz and then a polynomial sum of squares through \eqref{split}.
This proves the equivalence.
\end{proof}

\section{Polynomial multipliers and half-line completions}
The endpoint constraint on $J_{12}$ distinguishes polynomials from algebraic
functions. An odd algebraic function such as $v|v|$ can vanish quadratically
at zero without vanishing identically. This extra freedom suffices for the
source in Theorem~\ref{thm:obstruction}: the
\href{https://github.com/anubhav70747/q222-finite-sdp/tree/main/extracted-lifts/canonical-completion/halfline}{half-line construction}
gives an explicit rational quadratic selector on $v\geq0$, extended to
$v<0$ by reflection. Its four constant Gram blocks have orders $8,8,4,4$.
Adding source nonnegativity gives an affine pencil of total order $33$,
with $44$ auxiliaries. Its projection is closed.

A polynomial multiplier would instead require one fixed nonzero
nonnegative $q(v)$ such that every valid source admits a polynomial
$Z_c(v)$ satisfying
\[
\begin{pmatrix}q(P+K)&Z_c\\Z_c^T&q(N+K)\end{pmatrix}\PSD,\qquad
\begin{pmatrix}q(P+E)&q(K-E)-Z_c\\q(K-E)-Z_c^T&q(N+E)\end{pmatrix}\PSD
\quad(v\in\R).
\]
The first inequality forces $\deg Z_c\leq\deg q+2$. Univariate matrix
factorization would then give a fixed finite testing family for
$q(v)u^2D_uF_c$, with coefficient equations affine in the source and Gram
matrices. Since the multiplier is positive on a dense set, its certificate
would also imply validity of $F_c$ at every angle by continuity.
The
\href{https://github.com/anubhav70747/q222-finite-sdp/tree/main/extracted-lifts/canonical-completion/polynomial-obstruction}{fixed-multiplier theorem}
rules this out for every such $q$. It supplies a strictly valid rational
source family whose rescaled full-state slack tends to the Motzkin
polynomial plus a constant. Every fixed separable polynomial multiplier
leaves a square monomial in this limit, which can be cancelled from a
hypothetical polynomial sum of squares to obtain a contradiction.

A square-root factor on the half-line removes that limiting obstruction.
The quadratic half-line cone nevertheless excludes an all-positive valid
source, as shown by the exact dual in the
\href{https://github.com/anubhav70747/q222-finite-sdp/tree/main/extracted-lifts/canonical-completion/ramified-completion}{ramified completion}.
Adding selector terms $\sqrt v$ and $v^{3/2}$ repairs that example and gives
an affine cone of total order $49$. Whether this fixed enlarged family
contains every valid canonical source remains open.

\begin{thebibliography}{1}
\bibitem{Zalar} A. Zalar, \emph{Matrix Fej\'er--Riesz theorem with gaps},
Theorem 1.2, \href{https://arxiv.org/abs/1503.06034}{arXiv:1503.06034};
\href{https://zalara.github.io/Papers/MatrixFejerRieszArxivOriginal.pdf}{author's full text}.
\end{thebibliography}
\end{document}
